\chapter{Kiến thức nền tảng và công trình liên quan}
\label{ch:related}

Chương này (1) giới thiệu bài toán web agent và hạ tầng benchmark dùng
xuyên suốt khóa luận, sau đó khảo sát các công trình liên quan về (2) web
agent dựa trên LLM, (3) khám phá skill và workflow, (4) sinh dữ liệu huấn
luyện tự động cho agent, và (5) mô hình ngôn ngữ nhỏ trong vai trò agent,
trước khi (6) định vị khóa luận so với các hướng nghiên cứu này. Trong mỗi
chủ đề, các công trình được trình bày từ cũ đến mới.

\section{Kiến thức nền tảng}
\label{sec:rel-background}

\subsection{Bài toán web agent}

Một web agent nhận mục tiêu bằng ngôn ngữ tự nhiên và phải hoàn thành nó
bằng cách thao tác website như một người dùng: đọc trang rồi click, gõ và
điều hướng. Trong cách hình thức hóa dạng văn bản đã thành chuẩn, quan sát
là bản tuần tự hóa của accessibility tree, còn action space là một tập nhỏ
các thao tác giao diện có tham số (Mục~\ref{sec:method-graph} nêu bản cụ
thể hóa của chúng tôi). Chính sách của agent thường được xây trên khuôn mẫu
ReAct~\cite{yao2023react}: đan xen suy luận dạng văn bản với hành động, để
mỗi quyết định dựa trên một dòng suy nghĩ tường minh về trạng thái hiện
tại.

\subsection{Benchmark và môi trường}

Hạ tầng đánh giá phát triển theo hướng ngày càng sát thực tế.
WebShop~\cite{yao2022webshop} đưa ra một trang thương mại điện tử mô phỏng
với hơn một triệu sản phẩm thật và các chỉ dẫn mua sắm thu thập từ cộng
đồng; tuy nhiên, nó chỉ có một website đơn giản hóa với một thể loại task.
Mind2Web~\cite{deng2023mind2web} mở rộng lên hơn 2{.}000 task trên 137
website thật, nhưng đánh giá bằng cách phát lại các bản chụp website đã
lưu, nên agent được chấm theo độ khớp với chuỗi hành động ghi sẵn thay vì
theo việc thật sự đạt mục tiêu trong môi trường sống.
WebArena~\cite{zhou2024webarena} khép lại khoảng trống này bằng các bản sao
website đầy đủ chức năng, tự triển khai (một cửa hàng trực tuyến, trang
quản trị của nó, một instance GitLab và một diễn đàn kiểu Reddit), cùng
812 task dài hạn được chấm bằng \emph{bộ chấm tự động} trên trạng thái cuối.
Kết quả chủ đạo của nó định khung cho cả lĩnh vực: agent tốt nhất dựa trên
GPT-4 đạt 14{,}41\% tỷ lệ hoàn thành, so với 78{,}24\% của con người. Khóa
luận này dùng WebArena vừa làm môi trường khám phá vừa làm đích đánh giá
chính thức. Cuối cùng, hệ sinh thái
BrowserGym~\cite{dechezelles2025browsergym} chuẩn hóa giao diện quan sát và
hành động giữa các benchmark; runtime của chúng tôi xây trên nền này.

\subsection{Tinh chỉnh tiết kiệm tham số}

Fine-tune toàn bộ một mô hình 9B tham số vẫn rất tốn bộ nhớ.
LoRA~\cite{hu2022lora} thay vào đó đóng băng trọng số pretrain và chỉ huấn
luyện các ma trận cập nhật hạng thấp cộng thêm, giảm số tham số huấn luyện
đi nhiều bậc mà chất lượng hầu như không đổi; adapter sau đó có thể gộp
ngược vào trọng số gốc để serving. Toàn bộ fine-tuning trong khóa luận dùng
LoRA.

\section{Web agent dựa trên LLM}
\label{sec:rel-llm-agents}

Làn sóng đầu tiên của thời kỳ WebArena cải thiện \emph{kiến trúc agent}
quanh một mô hình lớn cố định. SteP~\cite{sodhi2024step} tổ chức một chồng
(stack) các chính sách LLM, tự động đẩy vào và lấy ra các chính sách con
(đăng nhập, tìm kiếm, điền form) như trạng thái điều khiển; nó vượt xa
baseline nguyên khối của WebArena. Tuy nhiên, thư viện chính sách của SteP
được thiết kế thủ công cho từng họ website, tái nhập chính công sức con
người mà lĩnh vực đang muốn loại bỏ.
AgentOccam~\cite{yang2024agentoccam} cho thấy chỉ cần tinh chỉnh cẩn thận
không gian quan sát và hành động, không cần bất kỳ bộ máy lập kế hoạch
nào, đã vượt nhiều hệ thống cầu kỳ đương thời trên WebArena, một minh
chứng tỉnh táo rằng phần lớn ``tiến bộ agent'' thực ra là kỹ thuật giao
diện. Tuy nhiên, cả hai hướng đều chung hai đặc điểm cấu trúc: chúng coi mô
hình là hàng hóa cố định (thường là hạng GPT-4, trả phí theo lời gọi), và
không để lại sản phẩm bền vững nào: những gì học được về một website
không sống sót qua episode. Với người vận hành cần agent rẻ, riêng tư và
cải thiện được, prompt khéo hơn trên mô hình độc quyền không giải quyết
ràng buộc; nó củng cố ràng buộc.

\section{Khám phá skill và workflow}
\label{sec:rel-skills}

Ý tưởng agent nên tích lũy các hành vi tái sử dụng có gốc rễ sâu. Khung
options~\cite{sutton1999options} hình thức hóa một skill thành bộ ba (tập
khởi đầu, chính sách, điều kiện kết thúc), còn STRIPS~\cite{fikes1971strips}
cho thấy các năng lực mô tả bằng precondition và effect hỗ trợ lập kế hoạch
ký hiệu, cả hai khung hình thức mà HSGE kế thừa trực tiếp
(Mục~\ref{sec:method-graph}).

Trong kỷ nguyên LLM, Voyager~\cite{wang2024voyager} minh họa việc học skill
suốt đời trong Minecraft: một thư viện ngày càng lớn các skill dạng mã thực
thi, mỗi skill được kiểm chứng bằng thực thi trước khi thu nhận. Tuy nhiên,
môi trường của nó có chân lý nền lập trình được (trạng thái game) mà
website không có, và skill của nó phục vụ runtime của một agent chứ không
phục vụ pipeline huấn luyện nào.
Synapse~\cite{zheng2024synapse} lưu nguyên trajectory làm mẫu (exemplar) và
truy xuất theo độ tương đồng để đưa vào prompt. Tuy nhiên, bộ nhớ
trajectory là phẳng: các mẫu không ghép được với nhau, và không gì phân
biệt một trajectory đã kiểm chứng với một trajectory ăn may.
Agent Workflow Memory~\cite{wang2024awm} rút ra các quy trình lặp lại
(workflow) từ trajectory quá khứ rồi tiêm vào prompt các lần sau, cải thiện
đáng kể tỷ lệ thành công trên WebArena. Tuy nhiên, workflow được lưu dưới
dạng văn bản phẳng, thu nhận không qua cổng kiểm chứng tường minh, và chỉ
được tiêu thụ như gợi ý lúc chạy; mô hình bên dưới không hề tốt lên.
Gần đây nhất, skill induction dạng lập trình~\cite{wang2025asi} nâng các
chuỗi hành động lặp lại thành hàm thực thi được, kiểm chứng ngay khi dùng,
ghép các hành động nguyên thủy thành những lời gọi ở mức cao hơn.
SkillWeaver~\cite{zheng2025skillweaver} đi xa hơn: agent tự khám phá, luyện
tập và tinh luyện skill thành các API gọi được thông qua khám phá lặp, cải
thiện đáng kể hiệu năng WebArena nhờ tái sử dụng và thậm chí chuyển giao
skill giữa các agent. Tuy nhiên, ở tất cả các hệ thống này, skill sinh ra
đều phục vụ việc tái sử dụng lúc suy luận bởi một mô hình lớn; không hệ nào
biến cấu trúc skill thành \emph{dữ liệu huấn luyện}.

HSGE đi theo hướng ngược lại: skill được phát hiện và kiểm chứng trên cùng
tinh thần, nhưng mục đích của chúng là được nâng thành các operator lập kế
hoạch mà các phép ghép chuỗi sẽ \emph{sinh ra} những task dài hạn để huấn
luyện một mô hình nhỏ.

\section{Sinh dữ liệu huấn luyện tự động cho agent}
\label{sec:rel-datagen}

Hướng nghiên cứu gần khóa luận nhất là sinh dữ liệu huấn luyện agent không
cần gán nhãn thủ công, phân nhánh theo nguồn dữ liệu.

\subsection{Từ tri thức gián tiếp}

Synatra~\cite{ou2024synatra} chuyển tri thức gián tiếp (các bài hướng
dẫn, tài liệu hướng dẫn thao tác trên mạng) thành bản mẫu tổng hợp ở quy mô lớn, tạo
ra 100k trajectory với chi phí thấp. Tuy nhiên, các bản mẫu được dựng
offline, không hề chạm vào website sống, nên không gì bảo đảm chúng thực
thi được trên môi trường thật.
AgentTrek~\cite{xu2024agenttrek} cải thiện độ bám thực tế bằng cách phát
lại các bài hướng dẫn web trong trình duyệt thật, có VLM kiểm chứng từng
lần phát lại. Tuy nhiên, độ phủ của nó bị chặn bởi những gì đã có bài hướng
dẫn: các quy trình chưa ai viết ra (bao gồm mọi thứ đặc thù của một
website riêng hoặc mới triển khai) nằm ngoài tầm với.

\subsection{Từ khám phá tự động}

Họ thứ hai, giống HSGE, để agent tự khám phá môi trường đích và chế tạo
task từ những gì tìm thấy.
NNetNav~\cite{murty2024nnetnav} khám phá website rồi gán nhãn hồi tố các
episode tương tác bằng chỉ dẫn, cắt tỉa những nhánh không chấp nhận nhãn
sub-task có nghĩa; fine-tune Llama-3.1-8B trên các bản mẫu này đạt trên
16\% trên WebArena. Tuy nhiên, khám phá là các rollout không cấu trúc:
những gì agent học được về năng lực của website bị vứt bỏ thay vì tích lũy,
mỗi episode lại bắt đầu từ con số không.
OS-Genesis~\cite{sun2024osgenesis} áp cùng ý tưởng tổng hợp ngược cho
trajectory GUI: tương tác trước, suy ra task sau. Tuy nhiên, task được suy
theo từng chuỗi tương tác, không có cơ chế ghép các năng lực phát hiện ở
các episode khác nhau.
PAE~\cite{zhou2024pae} khép vòng bằng học tăng cường: một bộ đề xuất task,
một chính sách hành động và một VLM đánh giá mà phán quyết trở thành phần
thưởng. Tuy nhiên, việc đề xuất task chỉ dựa trên ngữ cảnh cục bộ, và RL
trên phần thưởng nhị phân thưa vừa ngốn mẫu vừa nặng hạ tầng so với SFT.
InSTA~\cite{trabucco2025insta} đẩy công thức này lên bề rộng internet, sinh
task cho 150k website sống bằng LLM ở cả khâu sinh lẫn khâu chấm. Tuy
nhiên, bề rộng đánh đổi với chiều sâu: một lượt quét qua nhiều website
không thể khám phá các task nhiều bước, ghép nhiều năng lực, vốn đòi hỏi
tri thức bền vững về một website.
WebRL~\cite{qi2025webrl} tiến hóa một curriculum từ chính các thất bại của
agent, huấn luyện các mô hình mở trên WebArena-Lite với mức cải thiện mạnh,
còn WebAgent-R1~\cite{wei2025webagentr1} huấn luyện mô hình mở đầu--cuối
bằng học tăng cường nhiều lượt trên các rollout song song. Tuy nhiên, cả
hai đều giả định sẵn một phân phối task ban đầu để cải thiện, và vòng tự
tiến hóa của chúng cần bộ máy RL trực tuyến mà SFT tránh được. Một hướng bổ
trợ thay hẳn website sống: WebEvolver~\cite{fang2025webevolver} đồng tiến
hóa một world model với agent để sinh rollout tổng hợp mà không chạm môi
trường thật. Tuy nhiên, một trình mô phỏng học được thừa kế độ phủ của dữ
liệu đã huấn luyện nó, và trajectory của nó chỉ đáng tin ngang độ trung
thực của world model.

Gần khóa luận nhất, Go-Browse~\cite{gandhi2026gobrowse} đóng khung việc thu
thập dữ liệu thành tìm kiếm trên graph của một website: các trang đã thăm
trở thành các node thăm lại được, nên khám phá tái sử dụng thông tin giữa
các episode. Một mô hình 7B fine-tune trên dữ liệu thu được đạt 21{,}7\%
trên WebArena, state of the art cho mô hình dưới 10B tại thời điểm công
bố. Go-Browse xác nhận đúng tiền đề cốt lõi của khóa luận: khám phá
\emph{có cấu trúc} thắng rollout không cấu trúc về chất lượng dữ liệu. Tuy
nhiên, graph của nó là graph \emph{của các trang}, và task của nó vẫn là
các đề xuất bám trang được giải từ những trang đó. HSGE bổ sung một tầng
thứ hai mà Go-Browse không xây: các skill đã kiểm chứng được nâng thành
operator lập kế hoạch với precondition/effect chuẩn hóa. Các chuỗi đi qua
những operator đó sinh ra \emph{semantic task}: những mục tiêu trải dài
trên các năng lực phát hiện ở các episode khác nhau, điều mà không đề xuất
bám trang nào tạo ra được (Mục~\ref{sec:method-semantic}).

\section{Mô hình ngôn ngữ nhỏ trong vai trò agent}
\label{sec:rel-slm}

Luận điểm rằng mô hình nhỏ xứng đáng được đầu tư ngày càng được phát biểu
tường minh: một position paper gần đây~\cite{belcak2025slm} lập luận rằng
phần lớn các lời gọi agent đều hẹp và lặp lại, khiến SLM (rẻ hơn mười
đến ba mươi lần khi serving và fine-tune được trong vài giờ) là nền tảng
hợp lý về kinh tế cho các hệ thống agent.
Về thực nghiệm, AutoWebGLM~\cite{lai2024autowebglm} xây một web agent triển
khai được trên backbone 6B thông qua SFT theo curriculum và học theo sở
thích. Tuy nhiên, dữ liệu huấn luyện của nó khởi động từ các mẫu gán nhãn
thủ công và curriculum dựng tay. Các hệ thống dựa trên khám phá ở mục trước
cung cấp đúng nguồn dữ liệu còn thiếu: mô hình 8B của NNetNav và 7B của
Go-Browse cho thấy bản mẫu thu thập tự động dịch chuyển đáng kể mô hình nhỏ
trên WebArena, còn WebRL đẩy xa hơn bằng RL. Tuy nhiên, xuyên suốt các công
trình này, nút thắt nhất quán là pipeline dữ liệu chứ không phải bộ tối ưu
và đó chính là thành phần khóa luận thiết kế lại. Mô hình đích của
chúng tôi là Qwen3.5-9B, một mô hình trọng số mở gần đây mà họ mô hình được
tài liệu hóa trong báo cáo kỹ thuật Qwen3~\cite{yang2025qwen3}.

\section{Định vị khóa luận}
\label{sec:rel-positioning}

Bảng~\ref{tab:positioning} định vị HSGE theo năm trục: dữ liệu có đến từ
tương tác sống với website đích không; khám phá có được cấu trúc hóa bằng
một biểu diễn bền vững không; skill tái sử dụng có được kiểm chứng trước
khi dùng lại không; task dài hạn có được sinh bằng cách \emph{ghép} các
năng lực đã khám phá không; và pipeline có huấn luyện một mô hình nhỏ
không.

\begin{table}[H]
\centering
\caption[Định vị HSGE so với các hệ thống liên quan gần nhất]{Định vị HSGE
so với các hệ thống liên quan gần nhất. \CIRCLE{} = có,
\LEFTcircle{} = một phần, - = không.}
\label{tab:positioning}
\small
\begin{tabular}{@{}L{3.2cm} C{1.7cm} C{2.0cm} C{1.9cm} C{2.1cm} C{1.6cm}@{}}
\toprule
\textbf{Hệ thống} & \textbf{Tương tác sống} & \textbf{Khám phá có cấu trúc}
& \textbf{Skill kiểm chứng} & \textbf{Task dài hạn ghép chuỗi}
& \textbf{Huấn luyện SLM} \\
\midrule
Synapse~\cite{zheng2024synapse} & - & - & - & - & - \\
AWM~\cite{wang2024awm} & \CIRCLE & - & - & - & - \\
ASI~\cite{wang2025asi} & \CIRCLE & - & \CIRCLE & \LEFTcircle & - \\
SkillWeaver~\cite{zheng2025skillweaver} & \CIRCLE & - & \CIRCLE & \LEFTcircle & - \\
Synatra~\cite{ou2024synatra} & - & - & - & - & \CIRCLE \\
AgentTrek~\cite{xu2024agenttrek} & \CIRCLE & - & - & - & \CIRCLE \\
NNetNav~\cite{murty2024nnetnav} & \CIRCLE & - & - & - & \CIRCLE \\
OS-Genesis~\cite{sun2024osgenesis} & \CIRCLE & - & - & - & \CIRCLE \\
PAE~\cite{zhou2024pae} & \CIRCLE & - & - & - & \CIRCLE \\
WebRL~\cite{qi2025webrl} & \CIRCLE & \LEFTcircle & - & - & \CIRCLE \\
Go-Browse~\cite{gandhi2026gobrowse} & \CIRCLE & \CIRCLE & - & - & \CIRCLE \\
\textbf{HSGE (của chúng tôi)} & \CIRCLE & \CIRCLE & \CIRCLE & \CIRCLE & \CIRCLE \\
\bottomrule
\end{tabular}
\end{table}

Tóm lại, hướng kiến trúc agent chủ yếu cải thiện chiến lược prompting và điều phối quanh các mô hình lớn cố định. Phương pháp khám phá kỹ năng có khả năng tích lũy tri thức có cấu trúc, nhưng phần lớn chỉ sử dụng tri thức này trong giai đoạn suy luận. Trong khi đó, phương pháp sinh dữ liệu hướng đến huấn luyện mô hình nhỏ thường thực hiện khám phá theo cách thiếu cấu trúc hoặc chỉ tổ chức thông tin ở cấp độ trang. Theo hiểu biết của chúng tôi, hướng tiếp cận kết hợp các kỹ năng đã được kiểm chứng, ghép chúng thành các tác vụ ngữ nghĩa dài hạn và sử dụng những tác vụ này để huấn luyện mô hình ngôn ngữ nhỏ vẫn chưa được nghiên cứu đầy đủ. Đây chính là khoảng trống mà khóa luận hướng đến giải quyết.


\vspace{1em}
\noindent
Chương tiếp theo trình bày phương pháp lấp vào vị trí đó.
